\uuid{713r}
\titre{Changement de variables}
\auteur{hanani ; mzari}
\organisation{exo7}
\licence{CC BY-NC-SA 4.0}
\datecreate{2020-01-20}
\niveau{L1}
\module{Analyse}
\chapitre{Calcul d'intégrales}
\sousChapitre{Primitives diverses}
\difficulte{3}
\theme{Primitives, changement de variables}
\exo7qcmid{lille-2.primitives.027}

\contenu{
\begin{question}
\qclassification{127.11}
\qtags{motcle=changement de variables}
Le changement de variable $u=2+\cos x$ donne :
\begin{answers}
\good{$\displaystyle \int \sin x(2+\cos x)^5\mathrm{d}x=-\frac{(2+\cos x)^6}{6}+k$, où $k\in \Rr$.}
\bad{$\displaystyle \int (2+\cos x)^5\mathrm{d}x=\frac{((2+\cos x)^5)^6}{6}+k$, où $k\in \Rr$.}
\good{$\displaystyle \int \sin x\left(2+\cos x\right)^{-3}\mathrm{d}x=\frac{\left(2+\cos x\right)^{-2}}{2}+k$, où $k\in \Rr$.}
\bad{$\displaystyle \int \left(2+\cos x\right)^{-3}\mathrm{d}x=\frac{\left(2+\cos x\right)^{-2}}{2}+k$, où $k\in \Rr$.}
\end{answers}
\begin{explanations}
Avec $u=2+\cos x$, on a : $\mathrm{d}u=-\sin x\, \mathrm{d}x$. D'où
$$\int \sin x(2+\cos x)^5\mathrm{d}x=-\int u^5\mathrm{d}u=-\frac{u^6}{6}+k,\; k\in \Rr,$$
et
$$\int \sin x(2+\cos x)^{-3}\mathrm{d}x=-\int u^{-3}\mathrm{d}u=\frac{u^{-2}}{2}+k,\; k\in \Rr.$$
\end{explanations}
\end{question}
}
